\chapter{Триггеры}
\label{ch:ff}

До сих пор мы строили схемы \emph{без памяти}: сумматоры, компараторы,
дешифраторы. Их выход целиком определяется тем, что на входах \emph{сейчас}.
В охранной панели это уже даёт тревогу по условию «дверь открыта и охрана
включена», а при пороге шума ещё и «шум выше нормы». Но стоит злоумышленнику
закрыть дверь, как условие исчезнет, и сирена замолчит. Остаётся вопрос главы:
\emph{как схеме запомнить, что тревога была?} Ответ даст элемент, который
хранит один бит, --- триггер.

\begin{buildbox}
В Logra откройте вкладку \term{Обучение}, главу «Триггеры и регистры» книги
«из блоков», и пройдите урок «D-триггер: запомнить бит». Подайте на вход
несколько значений и посмотрите на вкладке \term{Диаграммы}: выход меняется
только в моменты фронта такта, а между фронтами хранит прошлое. Не читая
дальше, сформулируйте одной фразой, чем это отличается от схемы из
вентилей. Теория главы объяснит, как такой «кирпич» устроен изнутри.
\end{buildbox}

\section{Зачем схеме память}

Представьте кнопку «включить/выключить» на пульте: нажали --- свет
загорелся, отпустили --- продолжает гореть, нажали ещё раз --- погас. В момент,
когда кнопка отпущена, входы схемы одинаковы, а выход может быть и 0, и 1 ---
в зависимости от того, \emph{что было раньше}. Комбинационная схема так не
умеет. Нужна \term{память}.

Простейший элемент памяти хранит один бит и называется \term{триггером}
(от англ. \emph{trigger} --- «спусковой крючок»; в английской литературе ---
latch и flip-flop).

\begin{diybox}[Увидеть память]
\textbf{Задание.} \stars{1} Пройдите в Logra урок «Запомни один бит» (раздел
«Время и память» курса «Обучение»). \stars{2} Откройте пример
\code{edge-toggle} и сравните его поведение с примером \code{button}: что
происходит с выходом при повторном нажатии?\\
\textbf{В Logra:} вкладки Обучение, Код, Диаграммы.\\
\textbf{Ожидаемый результат:} у \code{button} выход повторяет вход,
у \code{edge-toggle} зависит ещё и от прошлого.\\
\textbf{Как проверить:} на Диаграммах найдите два момента, где входы одинаковы,
а выходы разные.
\end{diybox}

\section{Бистабильная ячейка}

Соединим два инвертора в кольцо. Если на выходе первого 1, то на выходе
второго 0, и этот 0 поддерживает 1 на выходе первого. Состояние
\emph{само себя поддерживает}. Возможно и противоположное устойчивое
состояние. У схемы два устойчивых состояния --- она \term{бистабильна} и
может хранить один бит.

\begin{figure}[H]
\centering
\begin{tikzpicture}
  \node[gNOT] (n1) at (0,0.9) {};
  \node[gNOT, rotate=180] (n2) at (0,-0.9) {};
  \draw[wire] (n1.output) -- ++(0.6,0) |- (n2.input); % fig-ok: feedback loop, entry segment 6 mm, rotated inverter
  \draw[wire] (n2.output) -- ++(-0.6,0) |- (n1.input);
  \node[font=\sffamily\small, text=FpgaRed] at (1.4,1.25) {1};
  \node[font=\sffamily\small, text=FpgaBlue] at (-1.75,-0.55) {0};
  \fill ($(n1.output)+(0.6,-0.9)$) circle (1.3pt);
  \draw[wire] ($(n1.output)+(0.6,-0.9)$) -- ++(0.6,0) node[right] {$Q$};
  \node[font=\sffamily\small, align=left, anchor=west] at (3.4,0.3) {Состояние поддерживает само себя,};
  \node[font=\sffamily\small, align=left, anchor=west] at (3.4,-0.3) {но изменить его пока нечем.};
\end{tikzpicture}
\caption{Два инвертора в кольце --- бистабильная ячейка}
\end{figure}

\begin{diybox}[Кольцо из двух инверторов]
\textbf{Задание.} \stars{2} В Конструкторе соедините два блока НЕ в кольцо и
выведите выход одного из них на светодиод. \stars{1} Если не хотите строить
петлю с нуля, прочитайте шаг «Память на один бит» в уроке «D-триггер: запомнить
бит» (глава «Триггеры и регистры»).\\
\textbf{Ожидаемый результат:} кольцо держит любое из двух состояний, но
записать в него новое значение нечем. Симулятор может показать неопределённое
значение \code{X}: у ячейки нет входа, задающего начальное состояние.\\
\textbf{Как проверить:} Диаграммы (выход не меняется сам) и RTL (две инверсии
по кругу).
\end{diybox}

\section{RS-триггер}

Чтобы записывать в ячейку новое значение, заменим инверторы на элементы
ИЛИ-НЕ и используем их вторые входы. Получится \term{RS-триггер}
(асинхронная RS-защёлка):
\begin{itemize}
  \item вход $S$ (Set --- «установить») записывает 1;
  \item вход $R$ (Reset --- «сбросить») записывает 0;
  \item если $S = R = 0$, триггер \textbf{хранит} прежнее значение.
\end{itemize}

\begin{figure}[H]
\centering
\begin{minipage}{0.48\textwidth}
\centering
\begin{tikzpicture}
  \node[gNOR] (g1) at (0,1.3) {};
  \node[gNOR] (g2) at (0,-1.3) {};
  \draw[wire] (g1.input 1) -- ++(-1.2,0) node[left] {$R$};
  \draw[wire] (g2.input 2) -- ++(-1.2,0) node[left] {$S$};
  \draw[wire] (g1.output) -- ++(1.4,0) node[right] {$Q$};
  \draw[wire] (g2.output) -- ++(1.4,0) node[right] {$\overline{Q}$};
  \draw[wire] ($(g1.output)+(0.5,0)$) |- (0,0.4) -| ($(g2.input 1)+(-0.45,0)$) -- (g2.input 1);
  \draw[wire] ($(g2.output)+(0.9,0)$) |- (0,-0.4) -| ($(g1.input 2)+(-0.85,0)$) -- (g1.input 2);
  \fill ($(g1.output)+(0.5,0)$) circle (1.3pt); \fill ($(g2.output)+(0.9,0)$) circle (1.3pt);
\end{tikzpicture}
\end{minipage}\hfill
\begin{minipage}{0.48\textwidth}
\centering
\begin{tabular}{cc|c|l}
\toprule
$S$ & $R$ & $Q$ & \textbf{Режим} \\
\midrule
0 & 0 & $Q$ & хранение \\
1 & 0 & 1 & установка \\
0 & 1 & 0 & сброс \\
1 & 1 & \textcolor{FpgaRed}{?} & \textcolor{FpgaRed}{запрещено} \\
\bottomrule
\end{tabular}
\end{minipage}
\caption{RS-триггер на двух элементах ИЛИ-НЕ и его таблица переходов}
\label{fig:rs}
\end{figure}

Разберём, почему он работает. Помним: выход ИЛИ-НЕ равен 0, если хотя бы один
вход равен 1.
\begin{itemize}
  \item $S=1$: нижний элемент выдаёт $\overline{Q}=0$. У верхнего оба входа
        равны нулю ($R=0$, $\overline{Q}=0$), поэтому $Q=1$.
  \item Теперь $S$ вернули в 0. На нижнем элементе $Q=1$, поэтому
        $\overline{Q}$ по-прежнему 0, а $Q$ по-прежнему 1: триггер
        \emph{запомнил} единицу.
  \item $S=R=1$ --- оба выхода становятся нулями (хотя должны быть
        противоположны), а при одновременном снятии сигналов неизвестно, в
        какое состояние триггер «упадёт». Поэтому это сочетание запрещено.
\end{itemize}

\begin{figure}[H]
\centering
\begin{tikztimingtable}[timing/slope=0.1, xscale=3.0, yscale=1.4, timing/font=\sffamily\small, timing/rowdist=3.2ex]
  $S$ & 1L 1H 10L \\
  $R$ & 5L 1H 6L \\
  $Q$ & 1L 4H 7L \\
\extracode
  \node[lbl, text=FpgaOrange] at (1.5,1.7) {уст.};
  \node[lbl, text=FpgaOrange] at (3.5,1.7) {хранение 1};
  \node[lbl, text=FpgaOrange] at (5.5,1.7) {сброс};
  \node[lbl, text=FpgaOrange] at (8.0,1.7) {хранение 0};
\end{tikztimingtable}
\caption{Работа RS-триггера: короткий импульс на $S$ или $R$ переключает
триггер, а без импульсов он хранит значение}
\end{figure}

\begin{infobox}[RS-триггер на элементах И-НЕ]
Точно так же триггер собирают из двух элементов И-НЕ. Но тогда входы
срабатывают от \emph{нуля}: $\overline{S}=0$ устанавливает, $\overline{R}=0$
сбрасывает, а запрещено сочетание $\overline{S}=\overline{R}=0$. Такой
триггер --- классическое средство подавления дребезга механического
переключателя: при первом же касании контакта он переключается и дальнейших
«прыжков» не замечает.
\end{infobox}

\begin{diybox}[RS-триггер]
\textbf{Задание.} \stars{1} Урок «SR-триггер: установить и сбросить»; в нём же
есть шаг «Запрещённая комбинация». \stars{2} Подайте $S=R=1$, затем снимите
оба сигнала и опишите, что произошло. \stars{3} Соберите «кнопка-выключатель
света» из двух кнопок: одна включает, другая выключает.\\
\textbf{В Logra:} Обучение, Конструктор, Диаграммы.\\
\textbf{Ожидаемый результат:} короткое нажатие $S$ включает, а свет остаётся
после отпускания; одновременное нажатие даёт неоднозначность.\\
\textbf{Как проверить:} Диаграммы: прогоните последовательность $S$, пауза,
$R$, пауза и сверьтесь с рис.~\ref{fig:rs}.
\end{diybox}

\subsection{RS-триггер на Verilog}
\label{sec:ff-sr-v}

Эту же схему можно описать на языке Verilog, не рисуя ничего мышью: каждый
элемент ИЛИ-НЕ превращается в одну строку \code{assign}. Выход каждого
элемента зависит от выхода другого --- получается та самая петля обратной
связи.

\vfile{\codepath/b7_flipflops/sr_latch.v}{code/b7\_flipflops/sr\_latch.v}

\textbf{Как читается.} Первая строка \code{assign} --- верхний элемент ИЛИ-НЕ:
$Q=\overline{R+\overline{Q}}$. Вторая --- нижний:
$\overline{Q}=\overline{S+Q}$. Ключевое слово \code{assign} описывает
\emph{провода и вентили}, а не последовательность действий, поэтому порядок
строк неважен, а то, что \code{q} и \code{qn} ссылаются друг на друга, как раз
и есть память. Тестбенч \code{tb\_sr\_latch.v} (рядом с модулем в папке
\code{code/b7\_flipflops}) делает четыре шага: $S=1$, затем $S=0$ (хранение
единицы), $R=1$, затем $R=0$ (хранение нуля). В начале нужно записать значение
(сначала $S=1$): пока триггер ничего не хранил, его выход неизвестен
(\code{X}), и это нормально.

\textbf{Как запустить в Logra.} Вкладка \term{Код}: вставьте модуль и нажмите
\term{Синтезировать} (Ctrl+Enter). Logra предупредит, что в модуле найдена
петля логики (\emph{logic loop}): для памяти из вентилей это ожидаемо, в
обычной комбинационной схеме такое предупреждение означало бы ошибку. На вкладке
\term{RTL} вы увидите два элемента ИЛИ-НЕ, замкнутых в кольцо; на
\term{Диаграммах} подайте последовательность $S$, пауза, $R$, пауза и сверьтесь
с рис.~\ref{fig:rs}.

\begin{diybox}[RS на Verilog: допиши модуль]
\textbf{Задание.} \stars{2} Опишите RS-триггер на элементах И-НЕ (входы
активны нулём: \code{sn} --- установка, \code{rn} --- сброс). \stars{3}
Выясните на Диаграммах, что произойдёт при \code{sn = rn = 0}, а затем
одновременно в \code{1} (запрещено на ИЛИ-НЕ, но не в этой схеме!). Чем
отличаются запрещённые комбинации у двух вариантов?\\
\textbf{Ожидаемый результат:} \code{sn = 0} даёт \code{q = 1}, \code{rn = 0} даёт
\code{q = 0}, при \code{sn = rn = 1} схема хранит.
\end{diybox}

\textbf{Образец ответа.}
\begin{vcode}
module sr_latch_n (
    input  sn,       // установка, активный уровень 0
    input  rn,       // сброс, активный уровень 0
    output q,
    output qn
);
    assign q  = ~(sn & qn);   // И-НЕ
    assign qn = ~(rn & q);    // И-НЕ
endmodule
\end{vcode}
Для \code{sn = rn = 0} оба выхода становятся единицами (запрещённая
комбинация этого варианта); после одновременного снятия сигналов состояние
неизвестно.

\section{D-защёлка}

У RS-триггера два неудобства: есть запрещённое сочетание входов, и он
реагирует на входы в любой момент. Добавим перед ним два элемента И, вход
данных $D$ и вход разрешения $C$ (clock, «синхронизация»). Получится
\term{D-защёлка} (D latch):
\begin{itemize}
  \item при $C=1$ защёлка \textbf{прозрачна}: выход повторяет вход, $Q = D$;
  \item при $C=0$ защёлка \textbf{хранит} значение, которое было на $D$ в момент
        перехода $C$ из 1 в 0.
\end{itemize}

\begin{figure}[H]
\centering
\begin{tikzpicture}
  \node[gAND] (a1) at (3,1.2) {};
  \node[gAND] (a2) at (3,-1.2) {};
  \node[draw=FpgaGreen, fill=green!8, thick, minimum width=18mm, minimum height=36mm, font=\sffamily\small, align=center] (rs) at (6.6,0) {RS-\\триггер};
  \node[font=\sffamily\footnotesize] at ($(rs.west)+(0.25,1.2)$) {R};
  \node[font=\sffamily\footnotesize] at ($(rs.west)+(0.25,-1.2)$) {S};
  \node[gNOT, anchor=output] (n) at ($(a1.input 1)+(-0.9,0)$) {};
  \coordinate (Dp) at ($(a2.input 2)+(-3.0,0)$);
  \coordinate (Dt) at ($(Dp)+(0.5,0)$);
  \draw[wire] (Dp) node[left] {$D$} -- (a2.input 2);
  \fill (Dt) circle (1.3pt);
  \draw[wire] (Dt) -- (Dt |- n.input) -- (n.input);
  \draw[wire] (n.output) -- (a1.input 1);
  \coordinate (cb1) at ($(a1.input 2)+(-0.4,0)$);
  \coordinate (cb2) at ($(a2.input 1)+(-0.4,0)$);
  \draw[wire, FpgaBlue] ($(Dp |- 0,0)$) node[left, text=FpgaDark] {$C$} -- (cb1 |- 0,0);
  \draw[wire, FpgaBlue] (cb1 |- 0,0) -- (cb1) -- (a1.input 2);
  \draw[wire, FpgaBlue] (cb1 |- 0,0) -- (cb2) -- (a2.input 1);
  \fill[FpgaBlue] (cb1 |- 0,0) circle (1.3pt);
  \draw[wire] (a1.output) -- ($(rs.west)+(0,1.2)$);
  \draw[wire] (a2.output) -- ($(rs.west)+(0,-1.2)$);
  \draw[wire] ($(rs.east)+(0,1.2)$) -- ++(0.8,0) node[right] {$Q$};
  \draw[wire] ($(rs.east)+(0,-1.2)$) -- ++(0.8,0) node[right] {$\overline{Q}$};
\end{tikzpicture}
\caption{D-защёлка: запрещённое сочетание $S=R=1$ теперь невозможно, так как
на $S$ и $R$ подаются $D$ и $\overline{D}$}
\end{figure}

\begin{figure}[H]
\centering
\begin{tikztimingtable}[timing/slope=0.08, xscale=1.6, yscale=1.4, timing/font=\sffamily\small, timing/rowdist=3.2ex]
  $C$ & 2L 4H 4L 4H 2L \\
  $D$ & 1L 2H 1L 4H 3L 1H 4L \\
  $Q$ & 2L 1H 1L 6H 1L 1H 4L \\
\extracode
  \begin{pgfonlayer}{background}
    \fill[green!10] (2,0.6) rectangle (6,-5.2);
    \fill[green!10] (10,0.6) rectangle (14,-5.2);
  \end{pgfonlayer}
  \node[lbl, text=FpgaGreen] at (4,1.7) {прозрачна};
  \node[lbl] at (8,1.7) {хранит};
  \node[lbl, text=FpgaGreen] at (12,1.7) {прозрачна};
\end{tikztimingtable}
\caption{D-защёлка при $C=1$ передаёт все изменения $D$ на выход, при $C=0$
«замораживает» выход}
\label{fig:dlatch_t}
\end{figure}

\begin{diybox}[D-защёлка]
\textbf{Задание.} \stars{2} В Конструкторе соберите защёлку по рис.~\ref{fig:dlatch_t}
из вентилей: НЕ, два И и RS-триггер на двух ИЛИ-НЕ. \stars{3} Подайте
изменяющийся $D$ и разрешение $C$ в виде длинного импульса.\\
\textbf{В Logra:} Конструктор, Диаграммы.\\
\textbf{Ожидаемый результат:} пока $C=1$, выход повторяет $D$ (включая все его
изменения), при $C=0$ значение замораживается.\\
\textbf{Как проверить:} Диаграммы; сравните с рис.~\ref{fig:dlatch_t}.
Сохраните схему: она понадобится в следующем разделе.
\end{diybox}

\subsection{Защёлка на Verilog: \code{always @*}}
\label{sec:ff-latch-v}

На Verilog защёлку описывают процессом \code{always @*}: это комбинационный
блок, который «просыпается» при любом изменении входов, которые в нём читаются.
Если в блоке не все ветки присваивают значение выходу, синтезатор делает вывод,
что при остальных условиях выход надо \emph{помнить}, и ставит защёлку.

\vfile{\codepath/b7_flipflops/d_latch.v}{code/b7\_flipflops/d\_latch.v}

\textbf{Как читается.} «Всегда, когда меняется \code{c} или \code{d}: если
\code{c} равно 1, то \code{q} равно \code{d}». Ветки \code{else} нет: при
\code{c = 0} выходу ничего не присваивается, то есть он хранит старое значение.
Так и получается память. Здесь латч задуман, и его присваивание делают
\emph{блокирующим} (\code{=}), как принято в комбинационных блоках.

\textbf{Как запустить в Logra.} Вкладка \term{Код}, Ctrl+Enter. Logra
предупредит о выведенной защёлке (в журнале Yosys: \emph{Latch inferred}). В учебном модуле это
подтверждение, что вы получили то, что хотели; в рабочем проекте такое
сообщение --- повод проверить, не забыли ли вы ветку \code{else}. На вкладке
\term{RTL} найдите блок защёлки, на \term{Диаграммах} подайте $D$ в виде
быстрых импульсов при длинном $C=1$ и посмотрите, как $Q$ повторяет $D$.

\begin{diybox}[Защёлка на Verilog: измени модуль]
\textbf{Задание.} \stars{1} Измените защёлку так, чтобы она была прозрачна
при $c=0$, а не при $c=1$. \stars{2} Уберите из модуля условие и запишите
\code{q = d} безусловно: во что превратилась схема и что теперь пишет
Logra?\\
\textbf{Ожидаемый результат:} предупреждение о латче исчезает, схема --- просто
провод от \code{d} к \code{q}.
\end{diybox}

\textbf{Образец ответа.}
\begin{vcode}
module d_latch_low (input c, d, output reg q);
    always @* begin
        if (!c) q = d;       // прозрачна при c = 0
    end
endmodule
\end{vcode}
Без условия получается \code{always @* q = d;}: это комбинационная схема без
памяти, а предупреждение о латче пропадает.

\section{D-триггер}

Прозрачность защёлки --- это и достоинство, и недостаток. Если выход защёлки
через комбинационную логику идёт обратно на её же вход (как в счётчике), то
пока $C=1$, сигнал будет бегать по кругу, и результат станет непредсказуемым.
Нужен элемент, который запоминает значение в \emph{один момент времени} ---
по \emph{фронту} тактового сигнала. Это \term{D-триггер} (D flip-flop),
главный элемент памяти всей современной цифровой техники.

Его можно собрать из двух D-защёлок, включённых последовательно, с
противоположными сигналами разрешения. Такая схема называется
\term{«ведущий -- ведомый»} (master--slave):

\begin{figure}[H]
\centering
\begin{tikzpicture}[lt/.style={draw=FpgaGreen, fill=green!8, thick, minimum width=24mm, minimum height=18mm, font=\sffamily\small, align=center}]
  \node[lt] (m) at (0,0) {\quad ведущая\\\quad защёлка};
  \node[lt] (s) at (4.2,0) {\quad ведомая\\\quad защёлка};
  \node[font=\sffamily\footnotesize] at ($(m.west)+(0.2,0.45)$) {D};
  \node[font=\sffamily\footnotesize] at ($(m.west)+(0.2,-0.45)$) {C};
  \node[font=\sffamily\footnotesize] at ($(s.west)+(0.2,0.45)$) {D};
  \node[font=\sffamily\footnotesize] at ($(s.west)+(0.2,-0.45)$) {C};
  \draw[wire] ($(m.west)+(-1.2,0.45)$) node[left] {$D$} -- ($(m.west)+(0,0.45)$);
  \draw[wire] ($(m.east)+(0,0.45)$) -- node[above, lbl] {$Q_m$} ($(s.west)+(0,0.45)$);
  \draw[wire] ($(s.east)+(0,0.45)$) -- ++(1,0) node[right] {$Q$};
  \node[gNOT, anchor=output] (n) at ($(m.west)+(-0.5,-0.45)$) {};
  \coordinate (ck) at ($(n.input)+(-0.5,0)$);
  \draw[wire, FpgaBlue] ($(ck)+(-0.4,0)$) node[left, text=FpgaDark] {CLK} -- (n.input);
  \draw[wire, FpgaBlue] (n.output) -- ($(m.west)+(0,-0.45)$);
  \draw[wire, FpgaBlue] (ck) |- (2.4,-1.9) |- ($(s.west)+(0,-0.45)$);
  \fill[FpgaBlue] (ck) circle (1.3pt);
  \node[font=\sffamily\small, align=left, anchor=north west] at (-2.4,-2.2) {CLK $=0$: ведущая прозрачна, ведомая хранит;\\
  CLK $=1$: ведущая хранит, ведомая прозрачна;\\ $\Rightarrow$ $Q$ меняется только в момент перехода CLK из 0 в 1.};
\end{tikzpicture}
\caption{D-триггер, построенный по схеме «ведущий -- ведомый»}
\label{fig:ms}
\end{figure}

На схемах D-триггер рисуют прямоугольником, где вход тактового сигнала
отмечен треугольником (знак «срабатывает по фронту»):

\begin{figure}[H]
\centering
\begin{tikzpicture}
  \draw[thick, fill=green!8] (0,0) rectangle (1.8,2.2);
  \node[font=\sffamily\small] at (0.3,1.7) {D}; \node[font=\sffamily\small] at (1.5,1.7) {Q};
  \node[font=\sffamily\small] at (1.5,0.5) {$\overline{\text{Q}}$};
  \draw[thick] (0,0.3) -- (0.25,0.5) -- (0,0.7);
  \draw[wire] (-0.7,1.7) node[left] {$D$} -- (0,1.7);
  \draw[wire] (-0.7,0.5) node[left] {CLK} -- (0,0.5);
  \draw[wire] (1.8,1.7) -- (2.5,1.7) node[right] {$Q$};
  \draw[wire] (1.8,0.5) -- (2.1,0.5); \draw[thick] (2.18,0.5) circle (0.08); \draw[wire] (2.26,0.5) -- (2.5,0.5) node[right] {$\overline{Q}$};
  \draw[wire] (0.9,2.2) -- (0.9,2.7) node[above] {$\overline{\text{SET}}$}; \draw[thick, fill=white] (0.9,2.28) circle (0.08);
  \draw[wire] (0.9,0) -- (0.9,-0.5) node[below] {$\overline{\text{RST}}$}; \draw[thick, fill=white] (0.9,-0.08) circle (0.08);
  \node[font=\sffamily\small, align=left, anchor=west] at (4.2,1.6) {треугольник --- срабатывает по фронту};
  \node[font=\sffamily\small, align=left, anchor=west] at (4.2,1.0) {кружок --- активный уровень 0};
  \node[font=\sffamily\small, align=left, anchor=west] at (4.2,0.4) {SET, RST --- асинхронные установка и сброс:};
  \node[font=\sffamily\small, align=left, anchor=west] at (4.2,-0.1) {срабатывают сразу, не дожидаясь фронта};
\end{tikzpicture}
\caption{Условное обозначение D-триггера}
\end{figure}

\begin{figure}[H]
\centering
\begin{tikztimingtable}[timing/slope=0.08, xscale=1.6, yscale=1.4, timing/font=\sffamily\small, timing/rowdist=3.2ex]
  CLK            & 2L 4H 4L 4H 2L \\
  $D$            & 1L 2H 1L 4H 3L 1H 4L \\
  $Q$ защёлки    & 2L 1H 1L 6H 1L 1H 4L \\
  $Q$ триггера   & 2L 8H 6L \\
\extracode
  \begin{pgfonlayer}{background}
    \draw[FpgaOrange, dashed] (2,1.3) -- (2,-7.2);
    \draw[FpgaOrange, dashed] (10,1.3) -- (10,-7.2);
  \end{pgfonlayer}
  \node[lbl, text=FpgaOrange] at (2,1.6) {фронт}; \node[lbl, text=FpgaOrange] at (10,1.6) {фронт};
\end{tikztimingtable}
\caption{Сравнение: защёлка (рис.~\ref{fig:dlatch_t}) пропускает все изменения
$D$, пока CLK $=1$, а триггер берёт значение $D$ \emph{только в момент фронта}}
\label{fig:latch_vs_ff}
\end{figure}

\begin{diybox}[Защёлка против триггера]
\textbf{Задание.} \stars{1} Урок «D-триггер: запомнить бит». \stars{2} Подайте
один и тот же сигнал $D$ на защёлку из прошлого раздела и на блок «D-триггер»,
а тактом сделайте длинный импульс. Найдите момент, где их выходы расходятся.
\stars{3} Подайте $D=\overline{Q}$ сначала на триггер, затем на защёлку.\\
\textbf{В Logra:} Конструктор, Диаграммы, RTL.\\
\textbf{Ожидаемый результат:} триггер делит частоту на 2, защёлка при $C=1$
уходит в быстрые колебания (в симуляторе это \code{X} или чередование).\\
\textbf{Как проверить:} Диаграммы; сравните с рис.~\ref{fig:latch_vs_ff}.
\end{diybox}

\subsection{Триггер на Verilog: \code{always @(posedge clk)}}
\label{sec:ff-dff-v}

D-триггер описывается процессом, который просыпается \emph{только по
переднему фронту} тактового сигнала. Сравните с защёлкой: меняется всего одна
строка, а поведение --- принципиально другое.

\vfile{\codepath/b7_flipflops/d_ff.v}{code/b7\_flipflops/d\_ff.v}

\textbf{Как читается.} «По каждому переднему фронту \code{clk} запомни в
\code{q} то, что сейчас на \code{d}». Присваивание здесь \emph{неблокирующее}
(\code{<=}): все триггеры схемы одновременно берут старые значения и
одновременно их обновляют. Правило, которое избавит от половины ошибок: в
\code{always @(posedge ...)} --- только \code{<=}, в \code{always @*} --- только
\code{=}. Зачем так, мы увидим в главе~\ref{ch:regs}, где бит «пролетит» весь
сдвиговый регистр за один такт, если нарушить правило.

Тестбенч \code{tb\_latch\_ff.v} подаёт один и тот же сигнал $D$ на защёлку и
на триггер:

\vfile{\codepath/b7_flipflops/tb_latch_ff.v}{code/b7\_flipflops/tb\_latch\_ff.v}

\textbf{Как читается тестбенч.} Задача \code{chk} ждёт 1~нс после изменения и
сравнивает оба выхода с ожидаемыми; оператор \code{!==} считает неизвестное
\code{X} отличием, а не совпадением. В конце тестбенч печатает \code{OK} или
число ошибок.

\textbf{Как запустить в Logra.} Вкладка \term{Код}: вставьте \code{d\_ff}
(и отдельно \code{d\_latch}) и нажмите Ctrl+Enter. Готовая пара «защёлка и
триггер в одном модуле» с общими входами --- модуль \code{dff\_vs\_latch} из
главы~\ref{ch:bridge}. На \term{Диаграммах} подайте на \code{clk} длинный
импульс, а внутри него меняйте \code{d}. Выход защёлки меняется вместе с \code{d}, выход триггера ---
только в момент фронта (рис.~\ref{fig:latch_vs_ff}).

\begin{diybox}[Триггер на Verilog: измени модуль]
\textbf{Задание.} \stars{1} Сделайте триггер, который срабатывает по
\emph{заднему} фронту (из 1 в 0). \stars{2} Добавьте к модулю второй выход
\code{qn}, всегда равный $\overline{Q}$. \stars{3} Соберите из двух таких
модулей триггер «ведущий -- ведомый» (рис.~\ref{fig:ms}) и сравните с
\code{d\_ff}.\\
\textbf{Ожидаемый результат:} выход меняется в момент спада \code{clk}.
\end{diybox}

\textbf{Образец ответа.}
\begin{vcode}
module d_ff_neg (
    input      clk,
    input      d,
    output reg q,
    output     qn
);
    always @(negedge clk) q <= d;   // negedge: задний фронт
    assign qn = ~q;
endmodule
\end{vcode}

\section{Временные параметры триггера}

Чтобы триггер надёжно запомнил значение, вход $D$ должен быть неизменным в
течение некоторого времени до фронта и после него:
\begin{itemize}
  \item $t_\text{su}$ (setup, \term{время предустановки}) --- сколько $D$ должен
        быть стабилен \emph{до} фронта;
  \item $t_\text{h}$ (hold, \term{время удержания}) --- сколько \emph{после} фронта;
  \item $t_\text{co}$ (clock-to-output) --- через сколько после фронта меняется $Q$.
\end{itemize}

\begin{figure}[H]
\centering
\begin{tikzpicture}[x=1cm, y=1cm]
  \draw[very thick, FpgaBlue] (0,2.2) -- (4,2.2) -- (4.1,3) -- (8,3);
  \node[font=\sffamily\small, anchor=east] at (-0.1,2.6) {CLK};
  \draw[very thick] (0,1.1) -- (1.5,1.1) -- (1.7,1.5) -- (6,1.5) -- (6.2,1.1) -- (8,1.1);
  \draw[very thick] (0,1.5) -- (1.5,1.5) -- (1.7,1.1) -- (6,1.1) -- (6.2,1.5) -- (8,1.5);
  \node[font=\sffamily\small, anchor=east] at (-0.1,1.3) {$D$};
  \node[font=\sffamily\footnotesize] at (3.9,1.3) {стабильно};
  \draw[very thick, FpgaGreen] (0,0) -- (4.8,0) -- (4.9,0.4) -- (8,0.4);
  \node[font=\sffamily\small, anchor=east] at (-0.1,0.2) {$Q$};
  \fill[FpgaRed, opacity=0.15] (2.6,0.9) rectangle (4.05,1.7);
  \fill[FpgaOrange, opacity=0.2] (4.05,0.9) rectangle (4.8,1.7);
  \draw[dashed, FpgaGray] (4.05,3.3) -- (4.05,-0.5);
  \draw[<->] (2.6,-0.35) -- (4.05,-0.35) node[midway, below, font=\sffamily\small] {$t_\text{su}$};
  \draw[<->] (4.05,-0.35) -- (4.8,-0.35) node[midway, below, font=\sffamily\small] {$t_\text{h}$};
  \draw[<->] (4.05,-1.1) -- (4.85,-1.1) node[midway, below, font=\sffamily\small] {$t_\text{co}$};
\end{tikzpicture}
\caption{Время предустановки, удержания и задержка выхода триггера}
\label{fig:setup_hold}
\end{figure}

\begin{warnbox}[Метастабильность]
Если $D$ меняется внутри окна $t_\text{su}$--$t_\text{h}$, триггер может на
некоторое время «зависнуть» в промежуточном состоянии --- ни 0, ни 1 --- и
лишь потом случайно упасть в одно из них. Это \term{метастабильность}. Она
неизбежно возникает, когда на триггер приходит сигнал, не связанный с
тактом (например, от кнопки). Защита --- \term{синхронизатор} из двух
триггеров подряд: к моменту следующего фронта первый триггер почти наверняка
успокоится.
\end{warnbox}

\begin{diybox}[Асинхронная кнопка]
\textbf{Задание.} \stars{1} Урок «Синхронизатор: привязать кнопку к такту».
\stars{2} На Диаграммах переместите нажатие кнопки на разные моменты такта и
объясните, на каком такте результат становится видимым. \stars{3} Добавьте
«Детектор фронта: один импульс на нажатие» (одноимённый урок) после
синхронизатора.\\
\textbf{В Logra:} Обучение, Конструктор, Диаграммы.\\
\textbf{Ожидаемый результат:} выход синхронизатора меняется только по фронту,
запаздывая на два такта; детектор даёт ровно один импульс на нажатие, сколько
бы тактов кнопка ни была зажата.\\
\textbf{Как проверить:} Диаграммы: удержание кнопки 10 тактов даёт
один импульс.
\end{diybox}

\subsection{Синхронизатор и детектор фронта на Verilog}
\label{sec:ff-sync-v}

Синхронизатор --- это ровно два триггера подряд, и на Verilog он занимает три
строки:

\vfile{\codepath/b7_flipflops/sync2.v}{code/b7\_flipflops/sync2.v}

\textbf{Как читается.} Оба присваивания стоят в одном блоке и оба
неблокирующие: за один фронт \code{ff1} получает значение входа, а \code{ff2}
--- \emph{старое} значение \code{ff1}. Поэтому значение проходит оба
триггера за два такта, а не за один. Метастабильность может случиться
только в \code{ff1}, и у неё есть целый такт, чтобы затихнуть до \code{ff2}.

Детектор фронта запоминает прошлое значение сигнала и сравнивает с
настоящим:

\vfile{\codepath/b7_flipflops/edge_det.v}{code/b7\_flipflops/edge\_det.v}

\textbf{Как читается.} \code{prev} --- значение \code{in} в прошлом такте.
Выражение \code{in \& \textasciitilde prev} истинно, когда сейчас 1, а раньше
был 0: один такт на каждое нажатие. Это \emph{комбинационный} выход,
поэтому его пишут через \code{assign}, а не в тактируемом блоке.

Цепочка «кнопка $\to$ синхронизатор $\to$ детектор фронта» собрана в одном
модуле:

\vfile{\codepath/b7_flipflops/button_pulse.v}{code/b7\_flipflops/button\_pulse.v}

Тестбенч \code{tb\_button\_pulse.v} нажимает кнопку \emph{не по фронту}
такта, держит её 10 тактов, отпускает и считает импульсы: должен получиться
ровно один. Он же проверяет, что \code{sync2} отвечает не за один такт, а за
два.

\textbf{Как запустить в Logra.} Вкладка \term{Код}: вставьте
\code{button\_pulse}, Ctrl+Enter, \term{RTL}: три триггера и один элемент
«И» с инверсией. На \term{Диаграммах} подайте на \code{btn} длинный
импульс на 10 тактов и посмотрите, что \code{pulse} длится один такт и
отстаёт от кнопки на два. Для платы кнопку подключают к выводу в
\term{Выводах}; не забудьте, что механическая кнопка ещё и дребезжит
(глава~\ref{ch:examples}).

\begin{diybox}[Детектор фронта: измени модуль]
\textbf{Задание.} \stars{1} Сделайте детектор \emph{заднего} фронта
(отпускание кнопки). \stars{2} Сделайте детектор \emph{любого} изменения
сигнала (оба фронта). \stars{3} Составьте цепочку: нажатие кнопки переключает
светодиод (синхронизатор, детектор фронта, T-триггер).\\
\textbf{Ожидаемый результат:} один импульс на отпускание; два импульса на одно
нажатие и отпускание; светодиод меняется только по нажатию.
\end{diybox}

\textbf{Образец ответа.}
\begin{vcode}
module edge_fall (input clk, in, output pulse);
    reg prev = 1'b0;
    always @(posedge clk) prev <= in;
    assign pulse = ~in & prev;       // сейчас 0, раньше 1
endmodule

module edge_any (input clk, in, output pulse);
    reg prev = 1'b0;
    always @(posedge clk) prev <= in;
    assign pulse = in ^ prev;        // изменилось в любую сторону
endmodule

module button_toggle (input clk, btn, output reg led);
    reg ff1 = 1'b0, ff2 = 1'b0, ff3 = 1'b0;
    initial led = 1'b0;
    always @(posedge clk) begin
        ff1 <= btn; ff2 <= ff1; ff3 <= ff2;
        if (ff2 & ~ff3) led <= ~led;   // импульс переключает T-триггер
    end
endmodule
\end{vcode}

\section{Вход разрешения}

Часто нужно, чтобы триггер обновлялся не на каждом фронте, а только по
команде. Для этого добавляют вход разрешения CE (clock enable). Устроен он
просто: перед входом $D$ ставится мультиплексор, который при CE $=0$ подаёт
на вход триггера его же выход --- триггер переписывает сам себя и не меняется.

\begin{figure}[H]
\centering
\begin{tikzpicture}
  \draw[thick, fill=orange!15] (0,0.9) -- (0.9,0.5) -- (0.9,-0.5) -- (0,-0.9) -- cycle;
  \node[font=\sffamily\footnotesize] at (0.45,0) {MUX};
  \node[font=\footnotesize, anchor=west] at (0.02,0.5) {0}; \node[font=\footnotesize, anchor=west] at (0.02,-0.5) {1};
  \draw[thick, fill=green!8] (1.8,-0.8) rectangle (3.2,0.8);
  \node[font=\sffamily\small] at (2.05,0) {D}; \node[font=\sffamily\small] at (2.95,0) {Q};
  \draw[thick] (2.3,-0.8) -- (2.5,-0.6) -- (2.7,-0.8);
  \draw[wire] (0.9,0) -- (1.8,0);
  \draw[wire] (3.2,0) -- (4.4,0) node[right] {$Q$};
  \draw[wire] (3.8,0) -- (3.8,1.5) -| (-0.6,0.5) -- (0,0.5); \fill (3.8,0) circle (1.3pt);
  \draw[wire] (-1.2,-0.5) node[left] {$D$} -- (0,-0.5);
  \draw[wire, FpgaBlue] (0.45,-0.7) -- (0.45,-1.4) node[below, text=FpgaDark] {CE};
  \draw[wire, FpgaBlue] (2.5,-0.8) -- (2.5,-1.4) node[below, text=FpgaDark] {CLK};
\end{tikzpicture}
\caption{D-триггер с входом разрешения: мультиплексор + обратная связь}
\label{fig:ce}
\end{figure}

\begin{diybox}[Разрешение записи]
\textbf{Задание.} \stars{1} Урок «Регистр с разрешением записи». \stars{2}
Урок «Регистр со сбросом»: добавьте синхронный сброс и проверьте, что при
одновременных сбросе и записи побеждает сброс. \stars{3} Соберите то же без
готового блока: мультиплексор плюс блок «D-триггер» по рис.~\ref{fig:ce}.\\
\textbf{В Logra:} Обучение, Конструктор, Диаграммы, RTL.\\
\textbf{Ожидаемый результат:} при $\text{CE}=0$ выход хранит значение, как бы
ни менялся $D$.\\
\textbf{Как проверить:} Диаграммы (изменяйте $D$ при $\text{CE}=0$); RTL
покажет тот же мультиплексор перед триггером.
\end{diybox}

\subsection{Разрешение и сброс на Verilog}
\label{sec:ff-en-rst-v}

На Verilog вход разрешения --- это просто \code{if} внутри тактируемого
блока. Мультиплексор с обратной связью (рис.~\ref{fig:ce}) синтезатор
построит сам.

\vfile{\codepath/b7_flipflops/dff_en.v}{code/b7\_flipflops/dff\_en.v}

\textbf{Как читается.} «По фронту: если \code{en} равно 1, запомнить \code{d}».
Ветки \code{else} нет, и это \emph{правильно}: в тактируемом блоке отсутствие
присваивания означает «\code{q} остаётся прежним», то есть триггер сам себя
переписывает. Запись \code{output reg q = 1'b0} задаёт начальное значение после
включения питания.

Тестбенч \code{tb\_dff\_en.v} проверяет три такта: \code{en=0}~--- значение
не меняется; \code{en=1}~--- записывается; снова \code{en=0}~--- хранится.
В Logra: вкладка \term{Код}, Ctrl+Enter, \term{RTL} (мультиплексор перед
триггером), \term{Диаграммы}: держите \code{en=0} и меняйте \code{d}, потом
подайте импульс на \code{en}.

Теперь \term{сброс}. Он бывает двух видов, и разница видна уже в
первой строке блока:

\vfile{\codepath/b7_flipflops/dff_rst_sync.v}{code/b7\_flipflops/dff\_rst\_sync.v}

\vfile{\codepath/b7_flipflops/dff_rst_async.v}{code/b7\_flipflops/dff\_rst\_async.v}

\textbf{Как читается.} У \emph{синхронного} сброса в списке чувствительности
только \code{posedge clk}: \code{rst} проверяется внутри блока, то есть
только в момент фронта, и стоит как обычный вход данных (триггер получает
лишний элемент на входе \code{d}). У \emph{асинхронного} в списке есть
\code{or posedge rst}: триггер обнуляется сразу, как только \code{rst}
поднялся, не дожидаясь такта. В ПЛИС у триггера есть отдельный вход
асинхронного сброса (как RST на условном обозначении выше), поэтому
лишней логики на \code{d} не прибавляется. Порядок веток \code{if (rst) \ldots\ else
\ldots} важен: сброс записывают первым, тогда он сильнее записи.

\vfile{\codepath/b7_flipflops/tb_dff_rst.v}{code/b7\_flipflops/tb\_dff\_rst.v}

Тестбенч подаёт сброс \emph{между} фронтами: синхронный триггер этого не
замечает до ближайшего фронта, асинхронный обнуляется сразу. В Logra
повторите на \term{Диаграммах}: соберите оба модуля, держите \code{d=1},
поднимите \code{rst} между фронтами и сравните моменты, когда выход упал в 0.

\begin{diybox}[Сброс и разрешение: допиши модуль]
\textbf{Задание.} \stars{1} Добавьте к \code{dff\_en} синхронный сброс, который
сильнее записи. \stars{2} Сделайте асинхронную \emph{установку} в 1
(\code{set}) вместо сброса. \stars{3} Сделайте асинхронный сброс с активным
уровнем 0 (вход \code{rst\_n}): в списке чувствительности понадобится
\code{negedge}.\\
\textbf{Ожидаемый результат:} при \code{rst = 1} и \code{en = 1} выход равен 0.
\end{diybox}

\textbf{Образец ответа.}
\begin{vcode}
module dff_en_rst (
    input      clk, rst, en, d,
    output reg q
);
    always @(posedge clk) begin
        if (rst)     q <= 1'b0;    // сначала сброс: он сильнее
        else if (en) q <= d;
    end
endmodule

module dff_aset (
    input      clk, set, d,
    output reg q
);
    always @(posedge clk or posedge set) begin
        if (set) q <= 1'b1;        // асинхронная установка в 1
        else     q <= d;
    end
endmodule

module dff_arst_n (
    input      clk, rst_n, d,
    output reg q
);
    always @(posedge clk or negedge rst_n) begin   // сброс по спаду rst_n
        if (!rst_n) q <= 1'b0;                      // активный уровень 0
        else        q <= d;
    end
endmodule
\end{vcode}
Запись в списке и проверка внутри блока должны быть согласованы: \code{negedge
rst\_n} в списке --- значит \code{if (!rst\_n)} внутри.

\section{T-триггер и JK-триггер}

\term{T-триггер} (toggle, «переключатель») меняет своё состояние на
противоположное по каждому фронту, если на входе $T=1$, и хранит его при
$T=0$. Его легко сделать из D-триггера и элемента «Исключающее ИЛИ»:
$D = T \oplus Q$.

\begin{figure}[H]
\centering
\begin{minipage}{0.5\textwidth}
\centering
\begin{tikzpicture}
  \node[gXOR] (x) at (0,0) {};
  \draw[thick, fill=green!8] (1.8,-0.8) rectangle (3.2,0.8);
  \node[font=\sffamily\small] at (2.05,0) {D}; \node[font=\sffamily\small] at (2.95,0) {Q};
  \draw[thick] (2.3,-0.8) -- (2.5,-0.6) -- (2.7,-0.8);
  \draw[wire] (x.output) -- (1.8,0);
  \draw[wire] (3.2,0) -- (4.1,0) node[right] {$Q$};
  \draw[wire] (3.65,0) -- (3.65,1.4) -| ($(x.input 1)+(-0.5,0)$) -- (x.input 1); \fill (3.65,0) circle (1.3pt);
  \draw[wire] (x.input 2) -- ++(-0.9,0) node[left] {$T$};
  \draw[wire, FpgaBlue] (2.5,-0.8) -- (2.5,-1.3) node[below, text=FpgaDark] {CLK};
\end{tikzpicture}
\end{minipage}\hfill
\begin{minipage}{0.45\textwidth}
\centering
\begin{tabular}{c|c|l}
\toprule
$T$ & $Q$ после фронта & \\
\midrule
0 & $Q$ & хранение \\
1 & $\overline{Q}$ & переключение \\
\bottomrule
\end{tabular}
\end{minipage}
\caption{T-триггер из D-триггера и элемента «Исключающее ИЛИ»}
\end{figure}

Если на T-триггер постоянно подавать $T=1$, он переключается на каждом
фронте, и частота на его выходе будет \emph{вдвое меньше} тактовой. Это
простейший \term{делитель частоты на 2}.

\begin{figure}[H]
\centering
\begin{tikztimingtable}[timing/slope=0.08, xscale=1.5, yscale=1.4, timing/font=\sffamily\small, timing/rowdist=3.2ex]
  CLK            & 8{1L 1H} \\
  $Q$ ($T=1$)    & 1L 2H 2L 2H 2L 2H 2L 2H 1L \\
\end{tikztimingtable}
\caption{T-триггер с $T=1$ делит частоту на 2}
\end{figure}

\textbf{T-триггер на Verilog.} Условие «если $T=1$, переключиться» ---
это один \code{if} и инверсия:

\vfile{\codepath/b7_flipflops/t_ff.v}{code/b7\_flipflops/t\_ff.v}

\textbf{Как читается.} По фронту: если сброс --- записать 0, иначе если
\code{t} равно 1 --- записать \emph{инверсию} самого себя (\code{~q}).
Запись \code{q <= ~q} читает старое значение \code{q} и записывает новое,
поэтому это не противоречие, а в точности «$D=T\oplus Q$». Если оставить
\code{t = 1}, получится делитель частоты на 2. Тестбенч \code{tb\_t\_ff.v}
проверяет, что при $t=0$ значение хранится, а при $t=1$ меняется на каждом
фронте. В Logra: вкладка \term{Код}, Ctrl+Enter; на \term{Диаграммах} держите
\code{t=1} и убедитесь, что период \code{q} вдвое больше периода \code{clk}
(как на рисунке выше); затем \code{t=0}~--- выход замирает.

\begin{diybox}[T- и JK-триггер на Verilog: допиши модуль]
\textbf{Задание.} \stars{2} Опишите JK-триггер: $J=1,K=0$ --- установка,
$J=0,K=1$ --- сброс, $J=K=1$ --- переключение, $J=K=0$ --- хранение.
\stars{3} Получите T-триггер из вашего JK-триггера соединением $J=K=T$.\\
\textbf{Ожидаемый результат:} при $J=K=1$ выход меняется на каждом фронте.
\end{diybox}

\textbf{Образец ответа.}
\begin{vcode}
module jk_ff (input clk, j, k, output reg q);
    always @(posedge clk) begin
        case ({j, k})                 // склейка двух битов в одно число
            2'b00: q <= q;            // хранение
            2'b01: q <= 1'b0;         // сброс
            2'b10: q <= 1'b1;         // установка
            2'b11: q <= ~q;           // переключение
        endcase
    end
endmodule
\end{vcode}

\term{JK-триггер} объединяет возможности RS- и T-триггеров: $J$ ---
установка, $K$ --- сброс, $J = K = 1$ --- переключение (запрещённого сочетания
нет). В старых схемах на микросхемах серии 155 JK-триггеры были очень
популярны. В современных микросхемах с программируемой логикой остаются
только D-триггеры, а остальные виды при необходимости получаются из них
добавлением логики на входе.

\begin{table}[H]
\centering
\caption{Сводка: виды триггеров}
\small
\begin{tabularx}{\textwidth}{l l X}
\toprule
\textbf{Тип} & \textbf{Когда срабатывает} & \textbf{Что делает} \\
\midrule
RS-защёлка & сразу & $S$ --- установить 1, $R$ --- сбросить в 0, $S=R=1$ запрещено \\
D-защёлка  & пока $C = 1$ & повторяет $D$, при $C = 0$ хранит \\
\textbf{D-триггер} & \textbf{по фронту} CLK & запоминает $D$ в момент фронта \\
T-триггер  & по фронту CLK & при $T=1$ меняет состояние на противоположное \\
JK-триггер & по фронту CLK & $J$ --- установка, $K$ --- сброс, $J=K=1$ --- переключение \\
\bottomrule
\end{tabularx}
\end{table}

\begin{ideabox}
Триггер хранит один бит. Главный элемент памяти в цифровой технике ---
\textbf{D-триггер}: он запоминает значение входа $D$ в момент \emph{переднего
фронта} тактового сигнала и хранит его до следующего фронта. Чтобы
триггер работал надёжно, вход не должен меняться в окне $t_\text{su}$--$t_\text{h}$
вокруг фронта.
\end{ideabox}

\begin{diybox}[T- и JK-триггеры; этап проекта]
\textbf{Задание.} \stars{1} Уроки «T-триггер: переключатель» и «JK-триггер».
\stars{2} Урок «Делитель частоты на 2» (глава «Счётчики»): убедитесь, что
выход вдвое медленнее такта. \stars{3} \textbf{Этап проекта «Охранная панель»:
запоминание тревоги.} Откройте проект «Сигнализация с запоминанием» (глава
«Проекты» книги «из блоков»): датчик \code{sensor} один раз дал 1, а
выход \code{y} остаётся 1, пока не придёт сброс \code{reset}. Датчиком
считайте сигнал «дверь открыта и охрана включена» из этапа~2, а сбросом ---
ключ охраны.\\
\textbf{В Logra:} Обучение (глава «Проекты»), Конструктор, Диаграммы, RTL.\\
\textbf{Ожидаемый результат:} дверь открыли на один такт, закрыли --- сирена
продолжает звучать; ключ сбрасывает тревогу, причём сброс сильнее датчика.\\
\textbf{Критерий приёмки:} все шесть проверок урока зелёные; на Диаграммах
видна «ступенька» выхода, которая не исчезает после закрытия двери; на RTL
один триггер и один элемент ИЛИ в цепи обратной связи.
\end{diybox}

\textbf{Этап проекта на Verilog.} Тот же узел в виде модуля:

\vfile{\codepath/b7_flipflops/alarm_keep.v}{code/b7\_flipflops/alarm\_keep.v}

\textbf{Как читается.} Две ветки в порядке приоритета: ключ сбрасывает,
иначе датчик устанавливает; если не пришло ни то, ни другое, \code{alarm}
хранит прежнее значение (ветки \code{else} нет). Тестбенч
\code{tb\_alarm\_keep.v} проходит семь шагов: тихо; дверь открыли; закрыли
(тревога помнится); ключ сбросил; ключ и датчик одновременно (ключ сильнее);
снова датчик. \textbf{Как запустить в Logra:} вкладка \term{Код}, вставьте
модуль, Ctrl+Enter, затем на \term{Диаграммах} подайте \code{sensor} на один
такт: \code{alarm} должен остаться единицей до импульса \code{key}.

\textbf{Допиши.} Добавьте вход \code{mute}, который выключает звук сирены, но
не сбрасывает тревогу. \emph{Образец:} к модулю добавляют порт
\code{output siren} и строку \code{assign siren = alarm \& \textasciitilde mute;}
--- сам триггер не меняется.

\section{Типичные ошибки}

\begin{warnbox}[Типичные ошибки]
\begin{enumerate}
  \item \textbf{«Защёлка --- это то же, что триггер».} Защёлка прозрачна
        всё время, пока $C=1$; триггер берёт значение только в момент фронта.
        Поставьте защёлку в петлю обратной связи (например, $D=\overline{Q}$) ---
        получится не делитель на 2, а генератор помех.
  \item \textbf{Снять оба сигнала $S=R=1$ одновременно.} Состояние после этого
        непредсказуемо: оба элемента одновременно «хотят» переключиться, и
        победит тот, кто окажется чуть быстрее.
  \item \textbf{Синхронизатор из двух триггеров ``навсегда решает
        проблему''.} Он лишь делает метастабильность на выходе чрезвычайно
        редкой, но не невозможной. И он годится для одного бита: шину из
        нескольких разрядов так передавать нельзя, разряды могут прийти в
        разных тактах.
  \item \textbf{Не задать начальное состояние.} После включения питания
        триггер находится в случайном состоянии. Если тревога «запомнена» из
        ничего, сирена будет выть сразу после включения. Нужен сброс.
  \item \textbf{Ждать, что $Q$ изменится в момент фронта.} Выход меняется
        через $t_\text{co}$ \emph{после} фронта, а входы собственного триггера
        «видят» старое значение. Именно поэтому цепочки триггеров сдвигаются,
        а не пролетают насквозь.
\end{enumerate}
\end{warnbox}

\subsection{Типичные ошибки на Verilog}

Пока триггеры были схемами, ошибки видны на картинке. Когда они
превращаются в код, появляются новые, и все они родом из одной особенности:
\emph{по тексту одинаково выглядят комбинационный блок и блок с памятью}.

\begin{warnbox}[Ошибки в описании памяти]
\begin{enumerate}
  \item \textbf{Случайная защёлка из-за неполного \code{if}.} Если в
        \code{always @*} выход присвоен не во всех ветках, синтезатор
        поставит защёлку, и Logra предупредит: \emph{Latch inferred}.
  \item \textbf{Нет \code{default}.} В \code{case} без ветки
        \code{default} (и без присваивания значения по умолчанию в начале блока)
        для неописанных комбинаций выход хранит прошлое значение: тоже защёлка.
  \item \textbf{Блокирующее \code{=} в тактируемом блоке и неблокирующее
        \code{<=} в комбинационном.} В \code{always @(posedge clk)} пишут
        \code{<=}; в \code{always @*} --- \code{=}. Нарушение правила даёт
        сдвиг, который «пролетает» за один такт (см.
        раздел~\ref{sec:regs-shift-v} в главе~\ref{ch:regs}).
  \item \textbf{Забыли значение по умолчанию после включения.} У \code{reg}
        до первой записи значение неизвестно (\code{X}); для счётчиков и
        автоматов нужен сброс или начальное значение при объявлении
        (\code{reg q = 0}).
  \item \textbf{Несогласованный сброс.} \code{posedge clk or posedge rst} в
        списке чувствительности и \code{if (!rst)} внутри блока: список
        и проверка должны соответствовать друг другу.
\end{enumerate}
\end{warnbox}

Первую ошибку удобно увидеть своими глазами:

\begin{vcode}
module oops (input sel, a, output reg y);
    always @* begin
        if (sel) y = a;      // нет else: y хранит значение, получилась защёлка!
    end
endmodule
\end{vcode}

Исправление --- дописать ветку или присвоить значение по умолчанию до
\code{if}:

\begin{vcode}
module oops_fixed (input sel, a, output reg y);
    always @* begin
        y = 1'b0;            // значение по умолчанию
        if (sel) y = a;
    end
endmodule
\end{vcode}

\section{Вопросы и задачи}
\label{sec:tasks-ff}

Ответы собраны в приложении~\ref{app:answers}.

\begin{taskbox}[Задачи к главе]
\begin{enumerate}
  \item \stars{1}~У RS-триггера на элементах ИЛИ-НЕ какая комбинация входов запрещена и почему?
  \item \stars{2}~Чем D-защёлка отличается от D-триггера? Нарисуйте реакцию обоих на один и тот же входной сигнал при тактовом импульсе.
  \item \stars{2}~Время установки 0,4~нс, задержка выхода триггера 0,5~нс, комбинационная логика между триггерами 3~нс. Каков минимальный период и какова максимальная частота?
  \item \stars{1}~Как из D-триггера сделать T-триггер? Что получится на выходе, если $T=1$ постоянно?
  \item \stars{1}~Что делает JK-триггер при $J=K=1$ по фронту такта?
  \item \stars{2}~Что такое метастабильность и в какой ситуации она возникает? Почему асинхронный сигнал пропускают через два триггера подряд?
  \item \stars{2}~\textbf{Найди ошибку.} Ученик хотел сделать делитель частоты на 2 и соединил выход $\overline{Q}$ D-защёлки с её входом $D$, а тактом сделал сигнал с частотой 1~МГц. Что получится и как исправить?
  \item \stars{2}~\textbf{Найди ошибку.} В «сигнализации с запоминанием» вход $D$ триггера соединили только с датчиком (без обратной связи через ИЛИ). Что произойдёт после закрытия двери?
  \item \stars{3}~\textbf{Проект.} Постройте в Конструкторе «сигнализацию с запоминанием» без готового урока: датчик \code{sensor}, ключ \code{key} (сброс), синхронный триггер. Критерий приёмки: на Диаграммах дверь открыта один такт, тревога держится, ключ сбрасывает, при одновременных \code{key} и \code{sensor} побеждает ключ.
  \item \stars{3}~\textbf{Проект.} Один импульс на нажатие: поставьте перед тревогой синхронизатор и детектор фронта. Критерий приёмки: кнопка зажата на 10 тактов, а счётчик событий (одно нажатие) увеличился на 1.
\end{enumerate}
\end{taskbox}

\section{Что дальше}

Мы научились хранить \emph{один} бит и знаем, когда ему можно доверять. Но
код замка, время тревоги и номер зоны --- это числа, а не отдельные биты.
Следующий вопрос: что получится, если поставить рядом несколько триггеров, и
как заставить их считать время для «писка» сирены? Об этом --- следующая
глава: регистры, сдвиг и счётчики.
